% !TeX root = ../main.tex

% Edit this title for each week's entry.
\section{Week 1}
\newlecture[Introduction/Syllabus][September 14, 2026]

\header{Principal and Interest}
\begin{definition}
    Principal: Amount invested

    Interest: "Rent" paid on investment
\end{definition}

In general, for
\begin{itemize}
    \item $A$: Principal
    \item $r$: interest rate per period
    \item $V$: Value received.
\end{itemize}
After 1 period, $V=A(1+r)$ is received

Interest over multiple periods depend on the compounding rules:
\begin{enumerate}
    \item \rb{Simple}: Money grows linearly, receiving $V=A(1+nr)$ over $n$ periods.
    \item \rb{Compound}: Money grows geometrically, receiving $V=A(1+r)^n$ over $n$ periods.
\end{enumerate}

The Rule of 72: Money invested at $r\%$ per year doubles in approximately $72/r$ years.


\begin{definition}
    \rb{Nominal Interest Rate}: annual quote, compounding interest is in fractions corresponding to the compounding period.
    \\For $m$ periods per year for $k$ periods, $V=(1+\frac{r}{m})^k$
\end{definition}

\begin{definition}
    \rb{Continuous Compound Interest}: when the compounding period becomes arbitrarily small, the value approaches $e^{rt}$.
\end{definition}

Relationship between Nominal and Effective Rate:
\[1+ r_{eff} = (1+ \frac{r_{nom}}{m})^m\]
\[\text{1 + Simple One-year Rate = the Explicit Compound}\]


\newlecture[Present Value and Internal Rate of Return][Sep 17, 2026]
Present value: comparing cash flows that occur at different times, while interest rates tells how to move cash flows in time. Use interest rates to 

\b{Ideal Bank}: applies the same rate of interest to deposits and loans, and has no transaction costs (real banks have higher rates for loans...).
Used for two ways:
\begin{itemize}
    \item \rb{Compounding}: when cash flow is moved \b{forward}, \dbb{multipled} by $(1+r)$ each period.
    \item \rb{Discounting}: when cash flow is moved \b{backward}, \dbb{divided} by $(1+r)$ each period.
\end{itemize}

\header{The Comparison Principle}
The ideal bank brings cash flow stream to the same time, then is compared.
\begin{itemize}
    \item When all cash flows are moved to \textit{present time}, it is called \rb{present value}.
    \item When all cash flows are moved to the \textit{same future timeline}, it is calle \rb{future value}.
\end{itemize}

\header{General Formula for Present Value}
Let $x = (x_0, x_1, x_2, \dots, x_n)$
\begin{align*}
    PV(x) &= x_0 + \frac{x_1}{(1+r)} + \frac{x_2}{(1+r)^2} + \cdots + \frac{x_n}{(1+r)^n} \\
    &= \boxed{\sum_{k=0}^{n}\frac{x_k}{(1+r)^k}}
\end{align*}
Note that it's usually assumed that $x_0$ is negative, for the initial investment.


\header{Main Theorem on Present Value}
Two cash flow streams are equivalent for an ideal bank iff \underline{they have the same present value}.

\header{Internal Rate of Return}
Captures the idea on the "return" on an investment.
\\The return on your investment is the return which would make the present value equal to 0.
\begin{example}
    Invest \$100 today, receive \$110 in one year. What is your return?
    \\Intuitively: 10\%, why?
    \[x(-100, 110)\]
    \[PV(x) = -100 + \frac{110}{1+0.1} = 0\]
\end{example}

\begin{definition}
    Let $x = (x_0, x_1, \dots, x_n)$ be a cash flow stream. Then the internal rate of return of this stream is the number $r$ satisfying:
    \[\boxed{\begin{gathered}
        PV(x) = 0 \\
        x_0 + \frac{x_1}{1+r} + \frac{x_2}{(1+r)^2} + \cdots + \frac{x_n}{(1+r)^n} = 0
    \end{gathered}}\]
    To solve for this $r$, its more convenient to set 
    \[c = \frac{1}{(1+r)}\]
    and solve
    \[x_0 + x_1c+ x_2 c^2 + \cdots + x_n c^n = 0\]
    with the fundamental theory of algebra.
\end{definition}

\begin{example}
    What is the internal rate of return of $(-2,1,1,1)$?
    \\Solve: $0 = -2 + c + c^2 + c^3$
    \[c = 0.81 \quad \Rightarrow \quad c = 0.81\]
\end{example}


\heading{Main Theorem on IRR}
Suppose the cash flow stream $x = (x_0, x_1, \dots, x_n)$ has $x_0 < 0$ and $x_k \geq 0$ where at least one term is strictly positive. Then, there is a unique positive root to the equation
\[x+0+x_1c + x_2 c^2 + \dots + x_n c^n = 0\]
Proof: plot $f(c) = x_0 + x_1c + x_2c^2 + \dots + x_n c^n$

\begin{center}
\begin{tikzpicture}
    \draw[->] (-0.3,0) -- (4.5,0) node[right] {$c$};
    \draw[->] (0,-1.5) -- (0,1.8) node[above] {$f(c)$};
    \draw[thick, blue, domain=0:4, smooth, variable=\x] plot ({\x}, {-1 + 0.15*\x*\x});
    \draw[dashed] (2.58,0) -- (2.58,-1.5);
    \filldraw (2.58,0) circle (1.5pt);
    \node[below] at (2.58,-0.05) {$c^*$};
    \filldraw (0,-1) circle (1.5pt);
    \node[left] at (0,-1) {$x_0$};
\end{tikzpicture}
\end{center}
Since $f(0) = x_0 < 0$ and $f'(c) = x_1 + 2x_2c + \dots + nx_nc^{n-1} \geq 0$ for $c \geq 0$ (as each $x_k \geq 0$, with at least one strictly positive), $f$ is nondecreasing and unbounded above. So $f$ crosses zero exactly once, at some $c^* > 0$.

\heading{Net Present Value (NPV)}
Computes all present values of \textit{all} cash flows associated with an investment. Includes the investment term $x_0$ as well.

\begin{example}
    Given two cash flow vectors
    \begin{enumerate}[label=(\arabic*)]
        \item After 1 year: $(-1,2)$
        \item After 2 years: $(-1,0, 3)$
        \\Using NPV solution (choose largest NPV):
        \[-1 + \frac{2}{1.1} = 0.82 \quad \text{vs} \quad -1 + \frac{0}{1.1} + \frac{3}{1.1^2} = 1.48\]
        NPV says choose option (2).
        \\Using IRR (choose greatest IRR):
        \[-1 + 2c = 0 \Rightarrow c = 0.5 \Rightarrow r = 1.0 \quad \text{vs} \quad -1+3c^2 = 0 \Rightarrow c = 0.577 \Rightarrow r = 0.7\]
        IRR says choose option (1).
    \end{enumerate}
\end{example}

\underline{IRR and NPV give conflicting recommendations.}
\\When scaled, \underline{IRR is scale-invariance, PV is not}. 
\\PV also \underline{distinguishes between single and repeated cash flows}, IRR does not.
\\In general, NPV is preferred.

\heading{Steps to Solving Cash Flow Problems}